% Transformative mathematical distillation prepared from a privately supplied
% transcription.  This file reconstructs the argument and formalization
% obligations without reproducing the source narrative.

\section{Section 8: selecting the physical branch of a double-angle estimate}

The sine and tangent double-angle inequalities control a function of
$2\Theta$.  Such a bound does not by itself distinguish a small angle from an
angle close to a right angle.  Section 8 resolves this ambiguity by attaching
the perturbed invariant subspace to a specified spectral component and then
showing that the continuation-selected component remains on the side
$\Theta<\pi/4$.

\subsection{Canonical spectral component}

Let $P$ reduce a bounded self-adjoint operator $A$, and suppose the spectrum on
$P\mathcal H$ lies below a cut $\alpha$ while the complementary spectrum lies
above $\alpha+\delta$.  Let $H$ be self-adjoint and off diagonal relative to
$P\mathcal H\oplus P^\perp\mathcal H$.  The canonical comparison subspace for
$A+H$ is not an arbitrary reducing subspace.  It is the spectral subspace on
the same side of the cut, equivalently the endpoint of the Riesz-projection
path continued from the original component.

If $Q$ denotes this selected projection, the crossed intersections
$P\mathcal H\cap Q^\perp\mathcal H$ and
$P^\perp\mathcal H\cap Q\mathcal H$ vanish.  The first follows because a vector
in the lower unperturbed component but the upper perturbed component would
have contradictory quadratic-form inequalities, while the off-diagonal
identity $PHP=0$ makes the two quadratic forms agree on that vector.  The
second intersection is excluded symmetrically.  Thus the selected pair is
acute and admits its canonical graph/direct-rotation description.

The modern Lean route represents the same selection by a common separating
contour along the affine path
\[
 A_t=A+tH,\qquad 0\le t\le1.
\]
The contour integral is identified with the genuine self-adjoint spectral
projection at every parameter.  Quantitative resolvent control makes the
projection path norm-continuous, finite subdivision gives locally close
orthogonal projections, and the local direct rotations compose to a unitary
transport between the endpoint spectral subspaces.

\subsection{The quarter-turn conclusion}

A selected endpoint that is a graph over the source component has an angular
coordinate $X$.  The projection gap satisfies
\[
 \|P-Q\|=\frac{\|X\|}{\sqrt{1+\|X\|^2}},
 \qquad
 \tan\theta_{\max}=\|X\|.
\]
Consequently $\|X\|<1$ is equivalent to
$\theta_{\max}<\pi/4$.  The continuation implementation currently derives
this contractive graph whenever its explicit common-contour Lipschitz
coefficient is below $\sqrt2/2$.

For the off-diagonal tangent theorem the historical argument is stronger: the
canonical branch stays below a quarter turn without imposing a small
perturbation norm.  Its scalar principal-plane calculation compares two ways
of writing the same perturbed block expectation.  If an angle crossed
$\pi/4$, the sign forced by the unperturbed gap would reverse the selected
spectral ordering.  The present repository has the Riccati and ordered-gap
pieces of this argument, but still needs their source-level assembly for the
unrestricted Theorem 8.1 hypotheses.

\subsection{Spectral repulsion}

Once the selected branch is known to occupy the lower half-line and its
orthogonal complement the upper half-line, exact block diagonalization gives
\[
 \operatorname{spec}(A+H)
 \subset(-\infty,\alpha]\cup[\alpha+\delta,\infty).
\]
This is genuine gap preservation, not merely an angle estimate.  Restricting
to the two endpoint subspaces yields pointwise separation of their spectra by
at least $\delta$.

The source proves finer statements.  In the direct-rotation coordinates, the
unperturbed upper compression is dominated by the perturbed upper compression
sandwiched by the cosine block, after shifting by the cut.  Finite-dimensional
min--max then gives ordered-eigenvalue repulsion, and a rearrangement/Ky Fan
argument yields a symmetric-gauge inequality with factors $\cos^2\theta_j$.
The quadratic-form algebra giving the compression inequality has now been
isolated and proved from an abstract rotated-block decomposition and the
sine/cosine Pythagorean identity.  Connecting that certificate to the
concrete direct-rotation blocks is still required.  The ordered-eigenvalue
and gauge refinements have not yet been assembled into stable source-facing
Lean declarations.

\subsection{The half-gap continuation theorem}

The second theorem begins with the general sine double-angle assumptions and
adds one of two strict half-gap conditions:
\[
 \|H\|<\frac\delta2
 \quad\hbox{or}\quad
 \|R\|<\frac\delta2.
\]
The original spectral component is assumed to lie in a central interval whose
endpoints are half a gap from the surrounding spectrum.  Along the reverse
path from $A+H$ to $A$, ordinary spectral enclosure leaves open intervals on
both sides of that central cluster.  The corresponding spectral projector is
therefore norm-continuous.  Starting at angle zero, the sine double-angle
estimate prevents a parameter value on the closed branch
$\theta\le\pi/4$ from reaching the boundary.  A connectedness argument then
keeps the entire path strictly below $\pi/4$.

For the residual alternative, the source invokes a replacement theorem that
changes the unused diagonal block of the perturbation without changing the
selected residual data or the perturbed restricted operators, and chooses the
replacement with norm equal to the residual norm.  This reduces the residual
case to the perturbation-norm case.  A source-complete formalization therefore
needs both the spectral-continuation construction from the half-gap hypothesis
and this replacement theorem.

\subsection{Current formal boundary}

The repository already proves the following reusable components without
admissions:

\begin{itemize}
\item common-contour Riesz projection construction and spectral-calculus
identification;
\item norm-continuity and quantitative Lipschitz control of the selected
projection path;
\item endpoint unitary transport;
\item strict quarter-acuteness from an explicit contour coefficient;
\item a unique contractive graph coordinate and its Riccati equation;
\item exact spectral union and gap exclusion once oriented branch placement is
supplied;
\item witness-level tangent estimates under ordered form bounds.
\end{itemize}

The Section 8 source package exposes these conclusions and records two narrow
bridges.  The first must construct the canonical contour and oriented branch
from the full off-diagonal hypotheses of Theorem 8.1.  The second must derive
the contour coefficient from either half-gap condition in Theorem 8.2, with a
separate Krein replacement argument for the residual alternative.  The
operator-order and symmetric-gauge repulsion refinements remain a subsequent
finite/infinite-dimensional promotion task.
